% ============================================================
%  Semelhança de Triângulos — Apostila Premium | PratiqueJá
%  Engine: XeLaTeX
% ============================================================
\documentclass[12pt]{article}
\usepackage{../../pratiqueja}
\usepackage{tikz}
\usetikzlibrary{calc}

% \hlm — destaque de resultado em modo-math (cor sem bold)
\newcommand{\hlm}[1]{{\color{darkblue}#1}}

\tikzset{
  tri/.style={draw=darkblue, line width=1.2pt, fill=accentblue!12!white},
  trj/.style={draw=darkblue, line width=1.2pt, fill=badgegreen!12!white},
  ra/.style={draw=darkblue, line width=0.8pt},
  hl/.style={draw=badgepurple!80!black, line width=1pt},
  cl/.style={font=\footnotesize\bfseries},
}

\setsubject{Semelhança de Triângulos}

\begin{document}
\pagenumbering{arabic}
\setstretch{1.2}
\color{bodytext}

\heroheader{Semelhança de Triângulos}%
           {Teorema de Tales, Critérios AA/LAL/LLL e Relações Métricas}%
           {Matemática · Intermediário}

\setlength{\columnsep}{18pt}
\setlength{\columnseprule}{0.5pt}
\def\columnseprulecolor{\color{exborder}}

\begin{multicols}{2}

%─────────────────────────────────────────────────────────────
\TW{Def}{badgepurple}{Definição}{O que é semelhança de triângulos?}

\regra{Dois triângulos são \textbf{semelhantes} quando têm ângulos
correspondentes \textbf{iguais} e lados correspondentes
\textbf{proporcionais}.}

\obs{$ABC \sim A'B'C'$ quando existe a \textbf{razão de semelhança} $k$:}
\formula{$\dfrac{AB}{A'B'} = \dfrac{BC}{B'C'} = \dfrac{AC}{A'C'} = k$}
\obs{$k = 1$ → triângulos \textbf{congruentes}. \quad $k \neq 1$ →
semelhantes de tamanhos diferentes.}

\nota{\textbf{Ordem dos vértices:} em $ABC \sim A'B'C'$, $A$ corresponde
a $A'$, $B$ a $B'$, $C$ a $C'$. Trocar a ordem altera quais lados são
correspondentes e invalida a proporção.}

%─────────────────────────────────────────────────────────────
\TW{TT}{darkblue}{Teorema de Tales}{Feixe de paralelas cortado por transversais}

\regra{Duas transversais cortando um feixe de \textbf{retas paralelas}
determinam segmentos \textbf{proporcionais}.}

\begin{center}\vspace{2pt}
\begin{tikzpicture}[scale=0.92]
\foreach \y in {2.2,0.8,-0.6}{
  \draw[line width=1pt, darkblue] (-0.5,\y)--(3.9,\y);}
\coordinate (L1) at (0.86,2.2); \coordinate (L2) at (0.44,0.8); \coordinate (L3) at (0.02,-0.6);
\coordinate (R1) at (2.54,2.2); \coordinate (R2) at (2.96,0.8); \coordinate (R3) at (3.38,-0.6);
\draw[line width=1pt, graycolor] (1.07,2.9)--(-0.19,-1.3);
\draw[line width=1pt, graycolor] (2.33,2.9)--(3.59,-1.3);
\foreach \p in {L1,L2,L3,R1,R2,R3}{\fill[vered] (\p) circle (1.6pt);}
\node[cl] at (0.95,1.5){$\alpha$};
\node[cl] at (0.53,0.1){$\beta$};
\node[cl] at (2.45,1.5){$\gamma$};
\node[cl] at (2.87,0.1){$\phi$};
\end{tikzpicture}
\end{center}

\formula{$\dfrac{\alpha}{\beta} = \dfrac{\gamma}{\phi}$}

\obs{A proporção mantém-se somando segmentos adjacentes:
$\dfrac{\alpha+\beta}{\beta} = \dfrac{\gamma+\phi}{\phi}$.}

\SH{Exemplos}
\exemp{%
  $\alpha=5,\ \beta=10,\ \phi=12$, achar $\gamma=x$:\\
  $\dfrac{5}{10} = \dfrac{x}{12} = \dfrac12 \Rightarrow x = \hlm{6}$\\[3pt]
  $\alpha=3,\ \beta=9,\ \gamma=2$, achar $\phi=x$:\\
  $\dfrac{3}{9} = \dfrac{2}{x} \Rightarrow 3x = 18 \Rightarrow x = \hlm{6}$%
}

%─────────────────────────────────────────────────────────────
\TW{AA}{badgegreen}{Critério AA}{Dois ângulos iguais garantem semelhança}

\regra{Se dois triângulos têm \textbf{dois ângulos} correspondentes
iguais, então são semelhantes (o terceiro também será, pois a soma é
$180°$).}

\begin{center}\vspace{2pt}
\begin{tikzpicture}[scale=0.95]
\coordinate (C) at (1.4,2.6); \coordinate (A) at (0,0); \coordinate (B) at (3.4,0);
\coordinate (D) at (0.7,1.3); \coordinate (E) at (2.4,1.3);
\draw[tri] (C)--(A)--(B)--cycle;
\draw[hl] (D)--(E);
\node[cl] at (1.4,2.85){$C$};
\node[cl] at (-0.22,-0.1){$A$};
\node[cl] at (3.6,-0.1){$B$};
\node[cl] at (0.42,1.38){$D$};
\node[cl] at (2.68,1.38){$E$};
\end{tikzpicture}
\end{center}

\obs{Com $DE \parallel AB$: $\angle CDE = \angle CAB$ e
$\angle CED = \angle CBA$, logo $\triangle CDE \sim \triangle CAB$:}
\formula{$\dfrac{CD}{CA} = \dfrac{CE}{CB} = \dfrac{DE}{AB}$}

%─────────────────────────────────────────────────────────────
\TW{LAL}{badgepurple}{Critério LAL}{Dois lados proporcionais e o ângulo entre eles}

\regra{Semelhantes quando \textbf{dois pares de lados} são proporcionais
e os \textbf{ângulos entre eles} são iguais. Se $\angle A = \angle D$ e
$\dfrac{AB}{DE} = \dfrac{AC}{DF}$, então $\triangle ABC \sim \triangle DEF$.}

\exemp{%
  $\triangle ABC$: $AB=4,\ AC=6,\ \angle A=50°$\\
  $\triangle DEF$: $DE=8,\ DF=12,\ \angle D=50°$\\[2pt]
  $\dfrac{AB}{DE} = \dfrac48 = \dfrac12$ \ e \ $\dfrac{AC}{DF} = \dfrac{6}{12} = \dfrac12$\\[2pt]
  ângulo incluído igual → \hlm{$\triangle ABC \sim \triangle DEF,\ k=2$}%
}

%─────────────────────────────────────────────────────────────
\TW{LLL}{darkblue}{Critério LLL}{Os três pares de lados proporcionais}

\regra{Semelhantes quando os \textbf{três pares} de lados
correspondentes são proporcionais — existe um único $k$ para todos.}

\exemp{%
  $\triangle ABC$: lados $3,\ 4,\ 6$ \qquad $\triangle DEF$: lados $6,\ 8,\ 12$\\[2pt]
  $\dfrac36 = \dfrac48 = \dfrac{6}{12} = \dfrac12$\\[2pt]
  todas as razões iguais → \hlm{$\triangle ABC \sim \triangle DEF,\ k=2$}%
}

%─────────────────────────────────────────────────────────────
\TW{Pr}{badgegreen}{Propriedades}{Perímetro, área, alturas e medianas}

\regra{Se a razão de semelhança é $k$, então:}
\formula{$\dfrac{P}{P'} = \dfrac{h}{h'} = \dfrac{m}{m'} = k$\\[4pt]
$\dfrac{A}{A'} = k^2$}
\obs{Perímetros, alturas e medianas crescem na razão $k$; a \textbf{área}
cresce com o \textbf{quadrado} $k^2$.}

\exemp{%
  $k=3$, área menor $=10\,\text{cm}^2$:\\
  $A' = 10\cdot3^2 = \hlm{90\,\text{cm}^2}$\\[3pt]
  $k=2$, perímetro menor $=12\,\text{cm}$:\\
  $P' = 12\cdot2 = \hlm{24\,\text{cm}}$%
}

%─────────────────────────────────────────────────────────────
\TW{Ret}{darkblue}{Triângulos Retângulos}{A altura sobre a hipotenusa gera 3 semelhantes}

\regra{Traçando a altura $h$ sobre a hipotenusa, formam-se \textbf{três
triângulos semelhantes}: $\triangle ABC \sim \triangle ACH \sim \triangle CBH$.}

\begin{center}\vspace{2pt}
\begin{tikzpicture}[scale=0.46]
\coordinate (A) at (0,0); \coordinate (B) at (10,0);
\coordinate (C) at (3.6,4.8); \coordinate (H) at (3.6,0);
\draw[tri] (A)--(B)--(C)--cycle;
\draw[hl] (C)--(H);
\draw[ra] (3.0,4.0)--(3.8,3.4)--(4.4,4.2);
\draw[ra] (3.6,0.7)--(4.3,0.7)--(4.3,0);
\node[cl] at (1.4,2.7){$c$};
\node[cl] at (7.2,2.9){$b$};
\node[cl] at (4.5,2.5){$h$};
\node[cl] at (1.8,-0.6){$m$};
\node[cl] at (6.8,-0.6){$n$};
\node[font=\scriptsize] at (3.6,-1.0){$H$};
\end{tikzpicture}
\end{center}

\obs{Sendo $m, n$ as projeções dos catetos sobre a hipotenusa
$a$ $(a=m+n)$:}
\formula{$h^2 = m\cdot n \qquad b^2 = a\cdot m \qquad c^2 = a\cdot n$}

\exemp{%
  Catetos $b=6,\ c=8$, achar $h$:\\
  $a = \sqrt{6^2+8^2} = 10$\\
  $m = \dfrac{36}{10} = 3{,}6$ \quad $n = \dfrac{64}{10} = 6{,}4$\\
  $h^2 = 3{,}6\cdot6{,}4 = 23{,}04 \Rightarrow h = \hlm{4{,}8}$\\
  {\footnotesize Verif.: $a\cdot h = 48 = b\cdot c$ \ok}%
}

\columnbreak

%─────────────────────────────────────────────────────────────
\TW{Prob}{badgegreen}{Problemas Contextualizados}{Semelhança em situações reais}

\regra{\textbf{Estratégia:} (1) identificar os triângulos semelhantes;
(2) montar a proporção entre lados correspondentes; (3) resolver.}

\SH{P1 — Altura de uma árvore}
\exemp{%
  Poste $2\,$m → sombra $3\,$m; árvore → sombra $12\,$m.
  Raios solares paralelos → semelhança (AA):\\
  $\dfrac{2}{h} = \dfrac{3}{12} \Rightarrow 3h = 24 \Rightarrow h = \hlm{8\,\text{m}}$%
}

\SH{P2 — Largura de um rio}
\exemp{%
  $DE \parallel AB$, $CD=10,\ DA=5,\ DE=8$. Achar $AB$:\\
  $CA = 10+5 = 15$\\
  $\dfrac{DE}{AB} = \dfrac{CD}{CA} \Rightarrow \dfrac{8}{AB} = \dfrac{10}{15} = \dfrac23$\\
  $2\cdot AB = 24 \Rightarrow AB = \hlm{12\,\text{m}}$%
}

\end{multicols}

\end{document}
